% Source document for the evaluation set: a typed teacher handout with planted traps.
% Traps: (1) wrong period T_0 = pi/omega_0 (must become \suspect, not be corrected),
% (2) crossed-out words (must become \struck), (3) on the simulated photos, the last
% line of page 2 is cut off (must become \missing).
\documentclass[11pt,a4paper]{article}
\usepackage[french]{babel}
\usepackage{amsmath,amssymb}
\usepackage[normalem]{ulem}
\usepackage[a4paper,margin=25mm]{geometry}
\pagestyle{empty}
\setlength{\parindent}{0pt}
\setlength{\parskip}{6pt}
\begin{document}
\textbf{\Large Chapitre 3 -- Oscillateur harmonique}

\textbf{Déf.} Un système de position $x(t)$ est un \emph{oscillateur harmonique} s'il vérifie
\[ \ddot{x} + \omega_0^2\, x = 0 \qquad (1) \]
où $\omega_0 > 0$ est la pulsation propre.

\textbf{Prop.} Les solutions de (1) sont les fonctions $x(t) = A\cos(\omega_0 t + \varphi)$, avec $A \geq 0$ et $\varphi \in \mathbb{R}$.

\textbf{Rq.} \sout{On admet que} Le mouvement est périodique, de période
\[ T_0 = \frac{\pi}{\omega_0}. \]

\textbf{Énergie.} Pour une masse $m$ reliée à un ressort de raideur $k$, on a $\omega_0 = \sqrt{k/m}$ et
\[ E = \frac{1}{2} m \dot{x}^2 + \frac{1}{2} k x^2 . \]

\newpage
\textbf{Ex.} Pour $m = 0{,}2$ kg et $k = 20$ N/m, on obtient $\omega_0 = 10$ rad/s.

\textbf{Algorithme (Euler symplectique)}

\begin{tabular}{l}
pour $n = 0$ à $N-1$ : \\
\quad $v_{n+1} \leftarrow v_n - h\,\omega_0^2\, x_n$ \\
\quad $x_{n+1} \leftarrow x_n + h\, v_{n+1}$ \\
\end{tabular}

\textbf{Exercice.} Montrer que $E$ est constante au cours du mouvement.

Indication : dériver $E$ par rapport au temps et utiliser (1).
\end{document}
